\chapter{Rayleigh Flow}

Rayleigh flow describes compressible, frictionless flow through a constant-area duct where heat addition or removal is the driving mechanism. Adding heat to a subsonic flow accelerates it toward $M = 1$ (the thermal choke point), while adding heat to a supersonic flow decelerates it toward $M = 1$. Removing heat has the opposite effect in both regimes. All flow properties at a given Mach number can be expressed as ratios to their values at the sonic reference state (denoted with an asterisk). Rayleigh flow is used to model combustion chambers, afterburners, and heat exchangers in high-speed propulsion systems. This module was recreated from the legacy toolkit codebase originally developed by Gowtham Sivaraman \cite{sivaraman2015}.

\section{Nomenclature}

\begin{center}
\begin{tabular}{@{}cl@{}}
\toprule
Symbol & Description \\
\midrule
$\gamma$ & Specific heat ratio \\
$M$ & Mach number \\
$\frac{T}{T^*}$ & Static to Sonic Temperature ratio \\
$\frac{P}{P^*}$ & Static to Sonic Pressure ratio \\
$\frac{\rho}{\rho^*}$ & Density to sonic density ratio \\
$\frac{T_0}{T_0^*}$ & Local stagnation to Sonic stagnation Temperature ratio \\
$\frac{P_0}{P_0^*}$ & Local stagnation to Sonic stagnation Pressure ratio \\
$\frac{U}{U^*}$ & Velocity to sonic velocity ratio \\
\bottomrule
\end{tabular}
\end{center}

\section{Governing Idea}

Every Rayleigh-flow property is a function of specific heat ratio $\gamma$ and Mach number $M$ alone. The solver therefore treats $M$ as the pivot variable: whichever quantity the user supplies is first inverted to recover $M$, after which every remaining property is evaluated directly from $M$ and $\gamma$. The problem is split into seven cases depending on which property is given.

\section{Calculation Cases}

Depending on the input parameter provided by the user, the Rayleigh flow properties are computed using one of the seven cases described below.

\subsection*{Case 1: Given $\gamma$ and $M$}

\textbf{Given:}
\begin{itemize}[nosep]
    \item Specific heat ratio $\gamma > 1$
    \item Mach number $M > 0$
\end{itemize}

\textbf{Calculation:} All properties follow directly from explicit analytical relations:
\begin{equation}
\frac{T}{T^*} = \frac{M^2(1+\gamma)^2}{(1+\gamma M^2)^2}
\end{equation}

\begin{equation}
\frac{P}{P^*} = \frac{1+\gamma}{1+\gamma M^2}
\end{equation}

\begin{equation}
\frac{\rho}{\rho^*} = \frac{1+\gamma M^2}{(1+\gamma)M^2}
\end{equation}

\begin{equation}
\frac{T_0}{T_0^*} = \frac{2(1+\gamma)M^2}{(1+\gamma M^2)^2}\left(1+\frac{\gamma-1}{2}M^2\right)
\end{equation}

\begin{equation}
\frac{P_0}{P_0^*} = \frac{1+\gamma}{1+\gamma M^2}\left(\frac{1+\frac{\gamma-1}{2}M^2}{\frac{\gamma+1}{2}}\right)^{\frac{\gamma}{\gamma-1}}
\end{equation}

\begin{equation}
\frac{U}{U^*} = \frac{(1+\gamma)M^2}{1+\gamma M^2}
\end{equation}


\subsection*{Case 2: Given $\gamma$ and $\frac{T}{T^*}$}

\textbf{Given:}
\begin{itemize}[nosep]
    \item Specific heat ratio $\gamma > 1$
    \item Temperature ratio $0 < \frac{T}{T^*} < \left(\frac{T}{T^*}\right)_{max}$, where the maximum temperature ratio is:
    \begin{equation}
    \left(\frac{T}{T^*}\right)_{max} = \frac{(1+\gamma)^2}{4\gamma}
    \end{equation}
\end{itemize}

\textbf{Calculation:} Two roots exist for the Mach number, split at the thermal choking point $M = 1/\sqrt{\gamma}$. Let:
\begin{equation}
AA = \sqrt{1 - 4\frac{T}{T^*}\gamma - 8\frac{T}{T^*}\gamma^2 - 4\frac{T}{T^*}\gamma^3 + 4\gamma + 6\gamma^2 + 4\gamma^3 + \gamma^4}
\end{equation}
Depending on the branch toggle selected (above or below $M = 1/\sqrt{\gamma}$):
\begin{itemize}
    \item \textbf{High-Mach branch toggle ($M > 1/\sqrt{\gamma}$):}
    \begin{equation}
    M = \frac{\sqrt{-2\frac{T}{T^*}\left(2\frac{T}{T^*}\gamma - 1 - 2\gamma - \gamma^2 - AA\right)}}{2\frac{T}{T^*}\gamma}
    \end{equation}
    \item \textbf{Low-Mach branch toggle ($M < 1/\sqrt{\gamma}$):}
    \begin{equation}
    M = \frac{\sqrt{-2\frac{T}{T^*}\left(2\frac{T}{T^*}\gamma - 1 - 2\gamma - \gamma^2 + AA\right)}}{2\frac{T}{T^*}\gamma}
    \end{equation}
\end{itemize}


\subsection*{Case 3: Given $\gamma$ and $\frac{P}{P^*}$}

\textbf{Given:}
\begin{itemize}[nosep]
    \item Specific heat ratio $\gamma > 1$
    \item Pressure ratio $0 < \frac{P}{P^*} < \left(\frac{P}{P^*}\right)_{max}$, where the maximum pressure ratio is:
    \begin{equation}
    \left(\frac{P}{P^*}\right)_{max} = 1+\gamma
    \end{equation}
\end{itemize}

\textbf{Calculation:} The Mach number is solved explicitly in closed form:
\begin{equation}
M = \frac{\sqrt{\frac{P}{P^*}\gamma\left(1+\gamma-\frac{P}{P^*}\right)}}{\gamma\times\frac{P}{P^*}}
\end{equation}


\subsection*{Case 4: Given $\gamma$ and $\frac{\rho}{\rho^*}$}

\textbf{Given:}
\begin{itemize}[nosep]
    \item Specific heat ratio $\gamma > 1$
    \item Density ratio $\frac{\rho}{\rho^*} > \left(\frac{\rho}{\rho^*}\right)_{min}$, where the minimum density ratio (at $M \to \infty$) is:
    \begin{equation}
    \left(\frac{\rho}{\rho^*}\right)_{min} = \frac{\gamma}{1+\gamma}
    \end{equation}
\end{itemize}

\textbf{Calculation:} The Mach number is solved explicitly in closed form:
\begin{equation}
M = \sqrt{\frac{1}{\frac{\rho}{\rho^*}+\gamma\frac{\rho}{\rho^*}-\gamma}}
\end{equation}


\subsection*{Case 5: Given $\gamma$ and $\frac{T_0}{T_0^*}$}

\textbf{Given:}
\begin{itemize}[nosep]
    \item Specific heat ratio $\gamma > 1$
    \item Stagnation temperature ratio $0 < \frac{T_0}{T_0^*} \le 1$. If the supersonic branch toggle is active, the ratio must be greater than the supersonic minimum limit:
    \begin{equation}
    \left(\frac{T_0}{T_0^*}\right)_{min,sup} = 1 - \frac{1}{\gamma^2}
    \end{equation}
\end{itemize}

\textbf{Calculation:} Two roots exist for the Mach number, split at the choking point $M = 1.0$. Let:
\begin{equation}
AA = \sqrt{-2\frac{T_0}{T_0^*}\gamma - \frac{T_0}{T_0^*}\gamma^2 + 1 + 2\gamma + \gamma^2 - \frac{T_0}{T_0^*}}
\end{equation}
The roots are solved explicitly based on the active branch:
\begin{itemize}
    \item \textbf{Subsonic branch toggle ($M < 1.0$):}
    \begin{equation}
    M_{sub} = \left|\frac{\sqrt{-\left(\frac{T_0}{T_0^*}\gamma^2+1-\gamma^2\right)\left(\frac{T_0}{T_0^*}\gamma - 1-\gamma+AA\right)}}{\frac{T_0}{T_0^*}\gamma^2+1-\gamma^2}\right|
    \end{equation}
    \item \textbf{Supersonic branch toggle ($M > 1.0$):}
    \begin{equation}
    M_{sup} = \left|\frac{\sqrt{-\left(\frac{T_0}{T_0^*}\gamma^2+1-\gamma^2\right)\left(\frac{T_0}{T_0^*}\gamma - 1-\gamma-AA\right)}}{\frac{T_0}{T_0^*}\gamma^2+1-\gamma^2}\right|
    \end{equation}
\end{itemize}


\subsection*{Case 6: Given $\gamma$ and $\frac{P_0}{P_0^*}$}

\textbf{Given:}
\begin{itemize}[nosep]
    \item Specific heat ratio $\gamma > 1$
    \item Stagnation pressure ratio $\frac{P_0}{P_0^*} \ge 1$
\end{itemize}

\textbf{Calculation:} No closed form exists; the Mach number $M$ is solved iteratively using the bisection method based on the selected flow branch toggle:
\begin{equation}
M = \frac{M_{start}+M_{end}}{2}
\end{equation}
\begin{equation}
f(M) = \frac{1+\gamma}{1+\gamma M^2}\left(\frac{1+\frac{\gamma-1}{2}M^2}{\frac{\gamma+1}{2}}\right)^{\frac{\gamma}{\gamma-1}} - \frac{P_0}{P_0^*} = 0
\end{equation}
The bisection solver updates the search brackets until $|f(M)| < 10^{-10}$ (with a limit of 200 iterations), using the following brackets based on the active branch:
\begin{itemize}[nosep]
    \item \textbf{Subsonic branch toggle:} Search interval $[M_{start}, M_{end}] = [10^{-9}, 1.0]$.
    \item \textbf{Supersonic branch toggle:} Search interval $[M_{start}, M_{end}] = [1.0, 100.0]$.
\end{itemize}


\subsection*{Case 7: Given $\gamma$ and $\frac{U}{U^*}$}

\textbf{Given:}
\begin{itemize}[nosep]
    \item Specific heat ratio $\gamma > 1$
    \item Velocity ratio $0 < \frac{U}{U^*} < \left(\frac{U}{U^*}\right)_{max}$, where the maximum velocity ratio (at $M \to \infty$) is:
    \begin{equation}
    \left(\frac{U}{U^*}\right)_{max} = \frac{1+\gamma}{\gamma}
    \end{equation}
\end{itemize}

\textbf{Calculation:} The Mach number is solved explicitly in closed form:
\begin{equation}
M = \sqrt{\frac{\frac{U}{U^*}}{1+\gamma-\gamma\frac{U}{U^*}}}
\end{equation}




\section{Program Flow and Architecture}

In Rayleigh flow, $M$ is the pivot variable. $P/P^*$, $\rho/\rho^*$, and $U/U^*$ are explicitly invertible; $T/T^*$ and $T_0/T_0^*$ have two roots (toggle selects branch); $P_0/P_0^*$ uses bisection.

\begin{center}
\begin{tikzpicture}[x=1cm,y=1cm,
    inp/.style={rectangle, draw=blue!70!black, thick, fill=blue!6, align=center,
                text width=4.5cm, minimum height=0.65cm, font=\footnotesize\sffamily},
    eng/.style={rectangle, draw=orange!80!black, thick, fill=orange!8, align=center,
                text width=3.4cm, font=\small\sffamily},
    outn/.style={rectangle, draw=green!60!black, thick, fill=green!6, align=center,
                text width=2.4cm, minimum height=0.65cm, font=\footnotesize\sffamily},
    note/.style={rectangle, draw=gray!60, dashed, fill=gray!5, align=center,
                 text width=2.0cm, minimum height=0.55cm, font=\scriptsize\sffamily},
    arr/.style={-{Stealth[scale=1.0]}, thick, draw=black!70},
    darr/.style={-{Stealth[scale=1.0]}, thick, dashed, draw=black!50},
    bus/.style={thick, draw=orange!70!black}
]
\node[inp] (i0) at (0, 0.0) {Case 1: $M$ (direct)};
\node[inp] (i1) at (0, 1.1) {Case 2: $T/T^*$ $\to$ explicit+toggle};
\node[inp] (i2) at (0, 2.2) {Case 3: $P/P^*$ $\to$ explicit};
\node[inp] (i3) at (0, 3.3) {Case 4: $\rho/\rho^*$ $\to$ explicit};
\node[inp] (i4) at (0, 4.4) {Case 5: $T_0/T_0^*$ $\to$ explicit+toggle};
\node[inp] (i5) at (0, 5.5) {Case 6: $P_0/P_0^*$ $\to$ bisection};
\node[inp] (i6) at (0, 6.6) {Case 7: $U/U^*$ $\to$ explicit};
\node[note] (t1) at (4.8, 1.1) {Sub/Sup\\toggle};
\node[note] (t2) at (4.8, 4.4) {Sub/Sup\\toggle};
\node[note] (t3) at (4.8, 5.5) {Sub/Sup\\toggle};
\node[eng, minimum height=7.2cm] (eng) at (7.6, 3.3)
    {\textbf{RayleighFlowEngine}\\[3pt]
    \texttt{fromMach(M,\,$\gamma$)}};
\node[outn] (o0) at (12.2, 0.0) {$M$};
\node[outn] (o1) at (12.2, 1.1) {$T/T^*$};
\node[outn] (o2) at (12.2, 2.2) {$P/P^*$};
\node[outn] (o3) at (12.2, 3.3) {$\rho/\rho^*$};
\node[outn] (o4) at (12.2, 4.4) {$T_0/T_0^*$};
\node[outn] (o5) at (12.2, 5.5) {$P_0/P_0^*$};
\node[outn] (o6) at (12.2, 6.6) {$U/U^*$};
\foreach \i in {i0,i1,i2,i3,i4,i5,i6} {\draw[arr] (\i.east) -- (eng.west |- \i);}
\draw[darr] (t1.east) -- (i1.east |- t1) -- (i1.east);
\draw[darr] (t2.east) -- (i4.east |- t2) -- (i4.east);
\draw[darr] (t3.east) -- (i5.east |- t3) -- (i5.east);
\coordinate (bx) at ($(eng.east)+(0.5,0)$);
\draw[bus] (eng.east) -- (bx);
\draw[bus] (bx |- o0) -- (bx |- o6);
\foreach \n in {o0,o1,o2,o3,o4,o5,o6} {\draw[arr] (bx |- \n) -- (\n.west);}
\end{tikzpicture}
\end{center}
